% Options for packages loaded elsewhere
\PassOptionsToPackage{unicode}{hyperref}
\PassOptionsToPackage{hyphens}{url}
\documentclass[
  11pt,
]{article}
\usepackage{xcolor}
\usepackage[margin=25mm]{geometry}
\usepackage{amsmath,amssymb}
\setcounter{secnumdepth}{-\maxdimen} % remove section numbering
\usepackage{iftex}
\ifPDFTeX
  \usepackage[T1]{fontenc}
  \usepackage[utf8]{inputenc}
  \usepackage{textcomp} % provide euro and other symbols
\else % if luatex or xetex
  \usepackage{unicode-math} % this also loads fontspec
  \defaultfontfeatures{Scale=MatchLowercase}
  \defaultfontfeatures[\rmfamily]{Ligatures=TeX,Scale=1}
\fi
\usepackage{lmodern}
\ifPDFTeX\else
  % xetex/luatex font selection
\fi
% Use upquote if available, for straight quotes in verbatim environments
\IfFileExists{upquote.sty}{\usepackage{upquote}}{}
\IfFileExists{microtype.sty}{% use microtype if available
  \usepackage[]{microtype}
  \UseMicrotypeSet[protrusion]{basicmath} % disable protrusion for tt fonts
}{}
\makeatletter
\@ifundefined{KOMAClassName}{% if non-KOMA class
  \IfFileExists{parskip.sty}{%
    \usepackage{parskip}
  }{% else
    \setlength{\parindent}{0pt}
    \setlength{\parskip}{6pt plus 2pt minus 1pt}}
}{% if KOMA class
  \KOMAoptions{parskip=half}}
\makeatother
\setlength{\emergencystretch}{3em} % prevent overfull lines
\providecommand{\tightlist}{%
  \setlength{\itemsep}{0pt}\setlength{\parskip}{0pt}}
\usepackage{bookmark}
\IfFileExists{xurl.sty}{\usepackage{xurl}}{} % add URL line breaks if available
\urlstyle{same}
\hypersetup{
  pdftitle={Finite observation windows for continuous powers of normal operators},
  pdfauthor={Anonymous},
  hidelinks,
  pdfcreator={LaTeX via pandoc}}

\title{Finite observation windows for continuous powers of normal
operators}
\author{Anonymous}
\date{5 November 2026 --- version 0.1.0-candidate}

\begin{document}
\maketitle

\textbf{Status:} Unrefereed mathematical manuscript, internally checked
within an AI-assisted research session. Historical priority remains
bounded by the literature search described below.

\textbf{Candidate identifier:}
\href{https://doi.org/10.5281/zenodo.23990540}{10.5281/zenodo.23990540}.
The
\href{https://github.com/ipitchford/finite-normal-observation-windows}{companion
repository} contains the sources and example checks. These identifiers
do not establish peer review or historical priority.

\subsection{Abstract}\label{abstract}

Let \(A\) be a bounded normal operator on a complex Hilbert space, and
let \(G\) be a finite or countable family of observation vectors. We
prove that if the continuous-power family
\(\{A^t g:g\in G,\ 0\leq t\leq T_0\}\) is a semi-continuous frame for
one finite \(T_0>0\), it is a semi-continuous frame for every finite
\(T>0\). The theorem requires neither a bounded inverse for \(A\),
reductivity, nor a Bessel assumption on \(G\). It therefore implies the
normal-reductive extension conjectured by Aldroubi, Huang and Petrosyan
and repeated in the 2026 dynamical-sampling survey. The proof obtains
locally bounded holomorphic observation maps directly from an integrated
upper frame bound. A reflected holomorphic energy and Montel's theorem
then propagate small observations from a shorter interval to the
original interval, contradicting its lower bound. Examples distinguish
injectivity from invertibility and show why delayed and infinite
observation windows require separate treatment. The conclusion concerns
existence of stable reconstruction on each window; it does not provide
uniform conditioning as the window shrinks.

\subsection{1. The observation-window
question}\label{the-observation-window-question}

Dynamical sampling studies reconstruction from measurements of a state
as it evolves. In the normal-operator model, the observation vectors at
time \(t\) have the form \(A^t g\). A frame inequality says that the
collected measurements determine the initial state with a bounded
reconstruction map. The lower inequality controls loss of information;
the upper inequality controls the total measurement energy.

Aldroubi, Huang and Petrosyan {[}1, Theorem 5.5{]} proved that a frame
on one finite interval beginning at zero remains a frame on every such
interval when \(A\) is invertible and self-adjoint. Their Conjecture 5.6
asks for a normal-reductive extension. Aldroubi, Cabrelli, Krishtal and
Molter {[}2, Theorem 2.13 and Conjecture 2{]} repeat this question. An
operator is \emph{reductive} if each of its closed invariant subspaces
is also invariant under its adjoint.

We establish a stronger assertion: normality suffices, and invertibility
is unnecessary. The distinction matters when spectral values approach
zero. In that case the logarithmic generator can be unbounded, and the
original observation family need not be Bessel even though its
observations integrated over time satisfy an upper frame bound.

The analytic-continuation strategy has a direct antecedent in the proof
of {[}1, Theorem 5.5{]}, which already uses Montel's theorem. The
central step here is a uniform holomorphic observation estimate that
works without a bounded logarithm or a Bessel observation family. This
supplies the lower bound uniformly over all unit initial states.
Analytic uniqueness for one fixed state would not suffice.

\subsection{2. Definitions and main
theorem}\label{definitions-and-main-theorem}

Let \(H\neq\{0\}\) be a separable complex Hilbert space. Inner products
are linear in the first variable. Let \(A\in\mathcal B(H)\) be normal,
meaning \(AA^*=A^*A\), with spectral measure \(E\). Fix a Borel argument

\[
\theta:\mathbb C\setminus\{0\}\longrightarrow[-\pi,\pi],
\qquad e^{i\theta(\lambda)}=\frac{\lambda}{|\lambda|}.
\]

For \(t>0\), set

\[
\lambda^t=\exp\!\bigl(t(\log|\lambda|+i\theta(\lambda))\bigr)
\quad(\lambda\neq0),\qquad 0^t=0,
\]

and define \(A^t\) by bounded Borel spectral calculus. Put \(A^0=I\).
The value at the single endpoint \(t=0\) does not affect any integral
below.

The proof also works for any other fixed bounded Borel choice of
argument. If its absolute value is bounded by \(B\), choose the
averaging interval in Section 3 sufficiently short that \(hB<2\pi\).
This prevents a zero of the averaging multiplier from entering the
spectral strip. No unbounded winding convention is asserted.

This choice specifies a spectral interpolation of the integer powers of
\(A\). The theorem is not a statement about every continuous-time
evolution having \(A\) as its time-one map.

Let \(G=(g_j)_{j\in J}\subset H\), where \(J\) is finite or countable.
Repetitions are allowed. Define the possibly infinite quadratic form

\[
Q_T(x):=\sum_{j\in J}\int_0^T
        |\langle x,A^t g_j\rangle|^2\,dt,
\qquad T>0.                                      \tag{2.1}
\]

The family is a \emph{semi-continuous frame} on \([0,T]\) if there exist
\(0<m_T\leq M_T<\infty\) such that

\[
m_T\|x\|^2\leq Q_T(x)\leq M_T\|x\|^2
\qquad(x\in H).                                  \tag{2.2}
\]

The unevolved family \(G\) is \emph{Bessel} if
\(\sum_j|\langle x,g_j\rangle|^2\leq B\|x\|^2\) for some finite \(B\).
We do not assume this property.

To make the reconstruction statement explicit, let \(\mathcal O_Tx\) be
the measurement function \(t\mapsto(\langle x,A^tg_j\rangle)_j\) in
\(L^2((0,T);\ell^2(J))\). Under (2.2) this is a bounded linear operator
with \[
\sqrt{m_T}\|x\|\leq\|\mathcal O_Tx\|\leq\sqrt{M_T}\|x\|.
\] Its range is closed and its inverse on that range is bounded. The
canonical reconstruction operator is
\((\mathcal O_T^*\mathcal O_T)^{-1}\mathcal O_T^*\). An additive
measurement error of norm at most \(\varepsilon\) therefore produces
reconstruction error at most \(\varepsilon/\sqrt{m_T}\). This is a
stability statement, not a computational prescription or a guarantee
that shorter observations are equally reliable.

The observed trajectory is \(U(t)x=(A^t)^*x\), measured against the
fixed sensors \(g_j\). We idealize the data as the complete
time-dependent output in this \(L^2\) norm; no finite selection of
sampling times, model-error bound or discretization-error bound is
supplied. Nor does the theorem make finitely many sensors sufficient in
infinite dimension. If \(G\) is finite, \(\mathcal O_T\) is
Hilbert--Schmidt, since its squared Hilbert--Schmidt norm is \[
\sum_{j\in J}\int_0^T\|A^tg_j\|^2\,dt<\infty.
\] A compact operator cannot be bounded below on an infinite-dimensional
Hilbert space. Thus finite \(G\) and the frame premise force
finite-dimensional \(H\). This familiar compactness limitation is
context, not a new contribution.

\subsubsection{Theorem 2.1. Invariance of finite observation
windows}\label{theorem-2.1.-invariance-of-finite-observation-windows}

Suppose \(A\) is bounded and normal and \(G\) is finite or countable. If
(2.2) holds for one \(T_0>0\), then it holds for every finite \(T>0\).

Consequently, the set of finite positive observation durations producing
a semi-continuous frame is either empty or all of \((0,\infty)\). No
assumption of invertibility, reductivity or the Bessel property of \(G\)
is needed.

The theorem asserts that each shorter window has some positive lower
frame bound. It does not assert a common lower bound for all durations,
or that the original lower bound can be retained.

Shorter windows can remain stable in this sense while becoming much more
sensitive to noise; Section 6.3 quantifies the distinction. The proof
has four steps. The lower bound first forces injectivity. The integrated
upper bound then controls observations at positive complex times, even
though the original sensor map may be unbounded. Reflection turns their
energy into a scalar analytic family. Finally, compactness rules out
unit states that become invisible on the shorter window but remain
visible on the original one. The averaging multiplier in Section 3
supplies the smoothing needed when the logarithm of the operator is
unbounded.

\subsection{3. Holomorphic observations from an integrated upper
bound}\label{holomorphic-observations-from-an-integrated-upper-bound}

We first isolate the estimate that permits an unbounded sensor analysis
map. Assume temporarily that \(\ker A=\{0\}\), and write

\[
b(\lambda)=\log|\lambda|+i\theta(\lambda),
\qquad w(\lambda)=\overline{b(\lambda)}.
\]

The logarithmic spectral values lie in the strip

\[
\operatorname{Re}w\leq a:=\log\|A\|,
\qquad |\operatorname{Im}w|\leq\pi.               \tag{3.1}
\]

For \(\operatorname{Re}z>0\), define

\[
U(z):=\int_{\sigma(A)\setminus\{0\}}
            e^{z w(\lambda)}\,dE(\lambda).       \tag{3.2}
\]

Then \(U(t)=(A^t)^*\) for real \(t>0\), and

\[
\|U(z)\|\leq
\max\{1,\|A\|^{\operatorname{Re}z}\}
e^{\pi|\operatorname{Im}z|}.                     \tag{3.3}
\]

The map \(z\mapsto U(z)\) is holomorphic in operator norm on the right
half-plane. On each compact subset, \(\operatorname{Re}z\) is bounded
below by a positive number. Exponential decay as
\(\operatorname{Re}w\to-\infty\) then bounds every multiplier
\(w^k e^{zw}\), including the derivatives required for operator-norm
holomorphy.

The use of the conjugate logarithm in (3.2) is deliberate. For a fixed
argument convention, \((A^*)^t\) need not equal \((A^t)^*\) on the
negative real spectral axis. For example, if \(\theta(-1)=-\pi\), the
two square-root expressions for \(A=-I\) are \(-iI\) and \(iI\).
Equation (3.2) avoids that ambiguity.

\subsubsection{Lemma 3.1. Positive-time holomorphic observation
bound}\label{lemma-3.1.-positive-time-holomorphic-observation-bound}

Suppose \(A\) is injective and normal and

\[
Q_{T_0}(x)\leq M\|x\|^2\qquad(x\in H).          \tag{3.4}
\]

For every \(x\in H\), the sequence

\[
F_x(z):=(\langle U(z)x,g_j\rangle)_{j\in J}      \tag{3.5}
\]

belongs to \(\ell^2(J)\) and defines a holomorphic function on
\(\operatorname{Re}z>0\). For each compact subset \(K\) of that
half-plane, there is a finite \(L_K\), independent of \(x\), such that

\[
\sup_{z\in K}\|F_x(z)\|_{\ell^2}\leq L_K\|x\|. \tag{3.6}
\]

\textbf{Proof.} Choose \(0<h\leq\min\{T_0,1\}\). Define

\[
r_h(w):=\int_0^h e^{sw}\,ds
       =\begin{cases}(e^{hw}-1)/w,&w\neq0,\\h,&w=0.
         \end{cases}                            \tag{3.7}
\]

The nonzero zeros of \(e^{hw}-1\) are \(2\pi i k/h\),
\(k\in\mathbb Z\setminus\{0\}\), and lie outside the strip (3.1). Thus
\(r_h\) has no zero there. Moreover,

\[
|r_h(w)|^{-1}\leq c_{h,a}(1+|w|)                \tag{3.8}
\]

for some finite constant \(c_{h,a}\). Indeed, when
\(\operatorname{Re}w\leq-1\),

\[
|e^{hw}-1|\geq1-e^{-h}.
\]

The rest of the strip is compact, if nonempty, and contains no zero of
\(r_h\).

It follows that

\[
d_z(w):=\frac{e^{zw}}{r_h(w)}                    \tag{3.9}
\]

is a bounded spectral multiplier, uniformly for \(z\) in each compact
subset of the right half-plane. Let \(D(z)\) be the corresponding
bounded operator. Set

\[
R_h:=\int_0^h U(s)\,ds.
\]

This operator-valued integral exists without norm continuity at zero.
For \(\varepsilon>0\), integrate the norm-continuous function on
\([\varepsilon,h]\). The omitted tail has norm at most
\(\varepsilon\sup_{0<s\leq h}\|U(s)\|\), so these integrals converge in
operator norm as \(\varepsilon\downarrow0\). Spectral calculus gives

\[
R_h=\int r_h(w(\lambda))\,dE(\lambda),
\qquad U(z)=R_hD(z).                             \tag{3.10}
\]

Enumerate \(J\), or use the finite family directly, and let

\[
C_Nx=(\langle x,g_j\rangle)_{j\leq N}.
\]

Each \(C_N\) is bounded. By (3.4),

\[
\int_0^h\|C_NU(s)x\|^2\,ds\leq M\|x\|^2.
\]

Hence Cauchy--Schwarz and (3.10) imply

\[
\begin{aligned}
\|C_NU(z)x\|
&=\left\|\int_0^h C_NU(s)D(z)x\,ds\right\|\\
&\leq\sqrt{hM}\,\|D(z)\|\,\|x\|.
\end{aligned}                                   \tag{3.11}
\]

The bound is independent of \(N\). Taking the supremum over finite
partial sums proves that (3.5) belongs to \(\ell^2\) and satisfies
(3.6).

For completeness, holomorphy also follows from these bounds without any
compactness assumption on the output space. On a closed disc strictly
inside the half-plane, finite-dimensional Cauchy estimates bound the
Taylor coefficient vectors of \(C_NU(z)x\) uniformly in \(N\). Their
coordinatewise limits belong to \(\ell^2\) with the same estimates. The
resulting vector-valued power series converges in \(\ell^2\) on each
smaller disc and agrees with every coordinate of \(F_x\). Thus \(F_x\)
is holomorphic. \(\square\)

The averaging multiplier resembles the bounded-generator admissibility
argument in {[}3, Theorem 3.9{]}. Here \(R_h\) need not have a bounded
inverse. The factor \(U(z)\) compensates for the growth of its
reciprocal at small spectral values, which is precisely what allows the
noninvertible case.

\subsection{4. Propagation of small observation
energy}\label{propagation-of-small-observation-energy}

The squared norm of a holomorphic vector-valued function is generally
not holomorphic. We use reflection in the real axis to obtain a scalar
holomorphic function with the required real-time values.

\subsubsection{Lemma 4.1. Propagation along positive real
times}\label{lemma-4.1.-propagation-along-positive-real-times}

Let \(Y\) be a complex Hilbert space, and let \(F_n\) be holomorphic
\(Y\)-valued functions on \(\operatorname{Re}z>0\), locally uniformly
bounded in norm. If

\[
\int_J\|F_n(t)\|^2\,dt\longrightarrow0           \tag{4.1}
\]

on a nondegenerate compact interval \(J\subset(0,\infty)\), then for
every compact interval \(K\subset(0,\infty)\),

\[
\sup_{t\in K}\|F_n(t)\|\longrightarrow0.         \tag{4.2}
\]

\textbf{Proof.} Define

\[
e_n(z):=\langle F_n(z),F_n(\overline z)\rangle_Y. \tag{4.3}
\]

This function is holomorphic: reflection in the variable and conjugate
linearity in the second inner-product slot cancel. The assertion also
follows by multiplying local power series. Cauchy--Schwarz gives local
uniform bounds for \(e_n\). For real \(t>0\),

\[
e_n(t)=\|F_n(t)\|^2\geq0.                       \tag{4.4}
\]

Montel's theorem implies that each subsequence has a further subsequence
converging uniformly on compact subsets to a scalar holomorphic function
\(e\). On \(J\), (4.1), uniform convergence and nonnegativity imply
\(e=0\). The identity theorem gives \(e\equiv0\) on the connected
half-plane. Every convergent subsequence therefore has the same zero
limit, so \(e_n\to0\) uniformly on compact subsets. Restricting to real
\(K\) and using (4.4) proves (4.2). \(\square\)

This lemma concerns a uniformly bounded analytic family. Applying the
identity theorem separately to each initial state would establish only a
uniqueness statement and would not yield the lower frame bound needed
here.

\subsection{5. Proof of the main
theorem}\label{proof-of-the-main-theorem}

Suppose

\[
m\|x\|^2\leq Q_{T_0}(x)\leq M\|x\|^2
\qquad(x\in H),\qquad m>0.                      \tag{5.1}
\]

\textbf{Injectivity.} If \(x\in\ker A\), normality gives
\(E(\{0\})x=x\). Thus \((A^t)^*x=0\) for every \(t>0\), and
\(Q_{T_0}(x)=0\). The lower bound forces \(x=0\). Lemma 3.1 therefore
applies.

\textbf{Upper bounds.} For \(T\leq T_0\), monotonicity of the integral
gives \(Q_T\leq Q_{T_0}\). For a longer interval, put
\(q=\lceil T/T_0\rceil\). The semigroup identity \(U(s+t)=U(s)U(t)\)
gives

\[
\begin{aligned}
Q_T(x)
&\leq\sum_{k=0}^{q-1}Q_{T_0}(U(kT_0)x)\\
&\leq M\sum_{k=0}^{q-1}\|A\|^{2kT_0}\,\|x\|^2,
\end{aligned}                                   \tag{5.2}
\]

where \(U(0)=I\). This is a finite upper bound.

\textbf{Lower bounds.} For \(T\geq T_0\), the lower bound in (5.1)
applies directly. Fix \(0<T<T_0\). If no positive lower bound existed,
there would be unit vectors \(x_n\) with

\[
Q_T(x_n)\longrightarrow0.                       \tag{5.3}
\]

By Lemma 3.1, the functions \(F_{x_n}\) are locally uniformly bounded
holomorphic \(\ell^2\)-valued functions. Their integrated squared norms
on \([T/3,2T/3]\) tend to zero. Lemma 4.1 gives

\[
\sup_{t\in[T,T_0]}\|F_{x_n}(t)\|\longrightarrow0.
\]

Consequently,

\[
\begin{aligned}
Q_{T_0}(x_n)
&=Q_T(x_n)+\int_T^{T_0}\|F_{x_n}(t)\|^2\,dt\\
&\longrightarrow0,
\end{aligned}
\]

contradicting \(Q_{T_0}(x_n)\geq m\). Thus \(Q_T\) has a positive lower
bound. This completes the proof of Theorem 2.1. \(\square\)

\subsection{6. Examples, conditioning and
boundaries}\label{examples-conditioning-and-boundaries}

\subsubsection{6.1. A noninvertible operator with non-Bessel
observations}\label{a-noninvertible-operator-with-non-bessel-observations}

On \(H=\ell^2(\mathbb N)\), let

\[
Ae_n=e^{-n}e_n,\qquad g_n=\sqrt{2n}\,e_n.
\]

The operator is injective and normal but has no bounded inverse. The
sensor family is not Bessel. Direct integration gives

\[
Q_T(x)=\sum_{n\geq1}(1-e^{-2nT})|x_n|^2.        \tag{6.1}
\]

Its optimal bounds are \(m_T=1-e^{-2T}\) and \(M_T=1\). Thus the
additional cases covered by the theorem are nonempty.

The unbounded instantaneous sensor energy is necessary for this
noninvertible case, not an incidental feature of the example. Suppose
instead that \(G\) has Bessel bound \(B\), and that the frame premise
holds. If \(A\) were not invertible, injectivity and normal spectral
calculus would provide unit vectors \(x_\varepsilon\) supported on
\(0<|\lambda|\leq\varepsilon\), for arbitrarily small \(\varepsilon<1\).
Then \[
Q_T(x_\varepsilon)\leq B\int_0^T\|U(t)x_\varepsilon\|^2\,dt
\leq B\int_0^T\varepsilon^{2t}\,dt
=B\frac{1-\varepsilon^{2T}}{2|\log\varepsilon|}\longrightarrow0,
\] contradicting the lower bound. Thus Bessel observations and the frame
premise force invertibility here. This elementary consequence explains
why allowing non-Bessel sensors genuinely enlarges the setting.

\subsubsection{6.2. Reductivity is
unnecessary}\label{reductivity-is-unnecessary}

Let \(A\) be multiplication by \(z\) on \(L^2\) of the unit circle, with
normalised arc measure, and take the Fourier orthonormal basis as \(G\).
The operator is unitary and normal. The closed span of the nonnegative
Fourier modes is invariant under \(A\) but not under \(A^*\), so \(A\)
is not reductive. Every \(A^t\), with the fixed Borel argument
convention, is unitary. Hence

\[
Q_T(x)=T\|x\|^2.
\]

\subsubsection{6.3. Shorter windows can be poorly
conditioned}\label{shorter-windows-can-be-poorly-conditioned}

Take \(A=\operatorname{diag}(1,e^{-1})\) on \(\mathbb C^2\), with one
observation vector \(g=(1,1)\). The observation Gramian is

\[
W_T=\begin{pmatrix}
T&1-e^{-T}\\
1-e^{-T}&(1-e^{-2T})/2
\end{pmatrix}.                                  \tag{6.2}
\]

The functions \(1\) and \(e^{-t}\) are linearly independent on every
positive-length interval, so \(W_T\) is positive definite. As
\(T\downarrow0\),

\[
\det W_T=\frac{T^4}{12}+O(T^5),\qquad
\operatorname{tr}W_T=2T+O(T^2),\qquad
\lambda_{\min}(W_T)=\frac{T^3}{24}+O(T^4).       \tag{6.3}
\]

The smallest possible norm of a left inverse for the observation map is
therefore asymptotic to \(\sqrt{24}\,T^{-3/2}\). Stable reconstruction
exists for each \(T\), while its sensitivity to measurement error
becomes arbitrarily large. These expansions follow directly from (6.2);
the accompanying \texttt{verify\_normal\_frames.py} checks the displayed
coefficients using exact rational series arithmetic. Its floating-point
integration checks are illustrations, not a verification of the abstract
theorem.

\subsubsection{6.4. Delaying observation can destroy the lower
bound}\label{delaying-observation-can-destroy-the-lower-bound}

For the example in (6.1), observations restricted to \([a,b]\), with
\(0<a<b<\infty\), give

\[
Q_{[a,b]}(x)=\sum_{n\geq1}
\bigl(e^{-2na}-e^{-2nb}\bigr)|x_n|^2.            \tag{6.4}
\]

The weights tend to zero. Testing on \(e_n\) shows that no positive
lower frame bound exists. Theorem 2.1 concerns shortening an interval
that still begins at zero, not moving the observation interval away from
zero. In fact there is a complete criterion, as a direct consequence of
the theorem and spectral calculus.

\textbf{Corollary 6.1 (delayed windows).} Under the hypotheses of
Theorem 2.1, fix \(a>0\) and \(T>0\). The observations on \([a,a+T]\)
form a frame if and only if \(A\) has a bounded inverse. In that case a
lower bound is \(m_T\|A^{-1}\|^{-2a}\), and an upper bound is
\(M_T\|A\|^{2a}\).

\textbf{Proof.} The semigroup identity gives
\(Q_{[a,a+T]}(x)=Q_T(U(a)x)\). Hence \[
m_T\|U(a)x\|^2\leq Q_{[a,a+T]}(x)\leq M_T\|U(a)x\|^2.
\] The delayed frame inequality holds exactly when \(U(a)\) is bounded
below. By normal spectral calculus its modulus is \(|A|^a\); this is
bounded below exactly when \(|A|\) is. A bounded normal operator bounded
below is invertible: its range is closed, and its orthogonal complement
is \(\ker A^*=\ker A=\{0\}\). The stated constants follow from the
spectral upper and lower bounds. \(\square\)

This corollary is an elementary consequence, not a separate priority
claim. It identifies the exact boundary behind the example: initial-time
measurements can compensate for rapidly decaying modes, but a positive
delay loses that compensation unless the spectrum stays away from zero.

\subsubsection{6.5. Finite and infinite windows are
different}\label{finite-and-infinite-windows-are-different}

If \(A=I\) and \(G\) is an orthonormal basis, then \(Q_T(x)=T\|x\|^2\)
for every finite \(T\), while the infinite-time upper bound fails for
every nonzero \(x\).

Conversely, let

\[
Ae_n=e^{-1/n}e_n,\qquad g_n=\sqrt{2/n}\,e_n.
\]

Then \(G\) is Bessel, \(A\) is bounded and invertible, and
\(Q_\infty(x)=\|x\|^2\). For finite \(T\), however,

\[
Q_T(x)=\sum_{n\geq1}(1-e^{-2T/n})|x_n|^2,
\]

whose weights tend to zero. No finite window has a positive lower bound.
Neither implication between finite-window and infinite-window frame
properties follows without further assumptions.

\subsection{7. Relation to observation theory and prior-art
boundary}\label{relation-to-observation-theory-and-prior-art-boundary}

The lower inequality in (2.2) is an exact initial-state observability
inequality. It should be distinguished from approximate observability,
which asserts uniqueness, and from final-state observability, whose
left-hand side controls an evolved state. The latter distinction matters
for parabolic control problems, where statements about minimum control
times concern different inequalities.

The dictionary in {[}3{]} makes the connection between dynamical frames
and control theory explicit. Its Theorem 3.11 states the self-adjoint
time-independence result. Engel and Nagel {[}4, Corollary VI.8.12{]}
give analytic-semigroup time independence for approximate
controllability. That theorem alone does not provide the uniform lower
norm estimate proved here.

Time independence is not a general property of strongly continuous
semigroups. Partington and Pott {[}5, Section 5, p.~30{]} explicitly
exhibit a left-shift semigroup whose exact-observation inequality holds
only for durations at least one. The following periodic transport
example gives a self-contained illustration with bounded observations.
Take

\[
H=L^2(\mathbb R/\mathbb Z),\qquad
T(t)f(s)=e^{-t}f(s+t\bmod1),\qquad
Cf=\mathbf1_{[0,a]}f,
\quad0<a<1.
\]

For \(T<1-a\), a nonzero function can avoid the observed arc throughout
the observation interval, so the lower bound fails. Over \([0,1]\),
every spatial point is observed for a total duration \(a\), and
\(e^{-2t}\geq e^{-2}\). Hence

\[
\int_0^1\|CT(t)f\|^2\,dt\geq ae^{-2}\|f\|^2.
\]

This semigroup is strongly continuous and exponentially stable; it is
not analytic. The example illustrates why the analytic structure used in
our proof matters. It is not claimed as a new counterexample or as a
refutation of a result in {[}5{]}.

The proof suggests a general analytic-semigroup principle whenever the
observation map has local holomorphic bounds away from zero. We do not
claim priority for that broader principle. Haak and Ouhabaz {[}6{]}
study exact observability with unbounded observations under additional
backward-forward conditioning and square-function assumptions. Those
additional conditions cannot be silently deduced from our single
finite-window frame premise; the distinction is particularly relevant to
the non-Bessel example.

The bounded literature search was refreshed on 5 November 2026. It
covered the original conjecture, exact observability and
controllability, admissible observations, analytic semigroups, bounded
generators, continuous-power frames and time-independence terminology.
The sources {[}1{]}--{[}6{]} were inspected in full or in their relevant
sections. Targeted Evidence Press catalogue/source searches and
OpenAI/TrureTuring corpus-title/path searches found no matching
published resolution; those corpus checks were not full-text reviews of
every paper. No earlier complete proof of the precise theorem presented
here was located within this boundary. This negative search result does
not establish historical priority, and older or unindexed control-theory
sources remain a limitation. A source-to-claim comparison accompanies
the evidence package.

The mathematical conclusion is independent of that priority question:
under the stated branch convention, a frame generated by a bounded
normal operator on one finite window beginning at zero remains a frame
on every such window. Determining useful lower constants and controlling
reconstruction error after substantial temporal compression remain
quantitative problems beyond the existence theorem.

\subsection{References}\label{references}

\begin{enumerate}
\def\labelenumi{\arabic{enumi}.}
\item
  Aldroubi, A., Huang, L. X., \& Petrosyan, A. (2019). Frames induced by
  the action of continuous powers of an operator. \emph{Journal of
  Mathematical Analysis and Applications, 478}(2), 1059--1084.
  https://doi.org/10.1016/j.jmaa.2019.05.066. Inspected preprint:
  https://arxiv.org/pdf/1801.10103.
\item
  Aldroubi, A., Cabrelli, C., Krishtal, I., \& Molter, U. (2026).
  Dynamical sampling: A survey. \emph{La Matematica, 5}, Article 37.
  https://doi.org/10.1007/s44007-026-00215-y.
\item
  Díaz Martín, R., Medri, I., \& Molter, U. (2021). Dynamical sampling:
  A view from control theory. In M. Hirn, S. Li, K. A. Okoudjou, S.
  Saliani, \& Ö. Yilmaz (Eds.), \emph{Excursions in harmonic analysis,
  volume 6} (pp.~269--295). Birkhäuser.
  https://doi.org/10.1007/978-3-030-69637-5\_15. Inspected preprint:
  https://arxiv.org/pdf/2003.01488.
\item
  Engel, K.-J., \& Nagel, R. (2000). \emph{One-parameter semigroups for
  linear evolution equations}. Springer. https://doi.org/10.1007/b97696.
  Inspected text:
  https://www.math.uni-tuebingen.de/de/forschung/agfa/members/engel-nagel\_one-parameter-semigroups.pdf.
\item
  Partington, J. R., \& Pott, S. (2005). Admissibility and exact
  observability of observation operators for semigroups. \emph{Irish
  Mathematical Society Bulletin, 55}, 19--39.
  https://www.maths.tcd.ie/pub/ims/bull55/R5502.pdf.
\item
  Haak, B. H., \& Ouhabaz, E. M. (2012). Exact observability, square
  functions and spectral theory. \emph{Journal of Functional Analysis,
  262}, 2903--2927. https://doi.org/10.1016/j.jfa.2012.01.007. Author
  manuscript:
  https://www.math.u-bordeaux.fr/\textasciitilde bhaak/recherche/HaakOuhabaz.pdf.
\end{enumerate}

\end{document}
